\section{Optimal Network Reconfiguration}

\subsection{Problem Form and Runtime Scope}

This chapter documents the current ONR mixed-integer model for radial feeder
reconfiguration.  The runtime problem is a LinDistFlow-style MILP with binary
branch-status variables, feeder-flow variables, squared-voltage variables, and
connectivity-flow variables, solved through the native branch-and-cut engine.

\subsection{Code Map}

The current implementation is concentrated in:
\begin{itemize}
  \item \texttt{src/analysis/topology\_analysis.cpp} for ONR model build, solve, and post-processing,
  \item \texttt{src/analysis/distribution\_power\_flow.cpp} for radial-feeder checks and distribution-side electrical support,
  \item \texttt{src/engine/branch\_and\_cut.cpp} and its helper files for MILP solution.
\end{itemize}

\subsection{Pseudocode Map}

The code-aligned workflow summary lives in two places:
\begin{itemize}
  \item Algorithm A13 in Section~\ref{sec:algorithms} for the application-level ONR solve flow,
  \item Section~\ref{sec:native-branch-and-cut} for the solver-kernel mechanics used underneath the ONR MILP.
\end{itemize}

\subsection{Current Status}

The ONR route is a live MILP application using the native branch-and-cut stack.
Its main current boundary is modeling scope: projected canonical AC components
are supported, but higher-level hybrid overlays such as VPP, microgrid, and
mobile-storage behavior are not yet injected directly into the ONR MILP.

\subsection{Scope of the Current ONR Module}

The ONR solver now exposes both \texttt{ACSystem} and
\texttt{HybridPowerSystem} overloads.  The \texttt{HybridPowerSystem} overload
first projects the rich model through \texttt{project\_to\_canonical\_models()},
so the current production path can see
\begin{itemize}
  \item direct AC buses and branches,
  \item direct AC loads and generators,
  \item direct static generators, renewable generators, PV systems, storage, and charging stations,
  \item projected transformers, switches, flexible loads, asymmetric loads, chargers, and optional three-phase reductions,
  \item still no automatic injection of \texttt{VPP}, \texttt{Microgrid}, or \texttt{MobileStorage} overlays into the ONR MILP.
\end{itemize}

\subsection{Decision Variables}

For an AC network with \(m=|\mathcal{E}_{ac}|\) branches and
\(n=|\mathcal{N}_{ac}|\) buses, the MILP uses the variable blocks
\begin{align}
x = \begin{bmatrix}
\alpha & P & Q & v & f & t
\end{bmatrix}^{\top},
\end{align}
where
\begin{align}
\alpha_b &\in \{0,1\} && \text{branch status}, \\
P_b,Q_b &\in \mathbb{R} && \text{branch active/reactive flow}, \\
v_i &\in \mathbb{R} && \text{squared bus voltage magnitude}, \\
f_b &\in \mathbb{R} && \text{single-commodity connectivity flow}, \\
t_b &\ge 0 && \text{auxiliary variable for } |P_b|.
\end{align}

\subsection{Net-Injection Model}

The ONR builder computes per-bus net injections directly from the AC component
tables as
\begin{align}
P_i^{net} &= \frac{P_i^{load} + P_i^{cs} - P_i^{gen} - P_i^{sg} - P_i^{ren} - P_i^{pv} - P_i^{stor}}{S_{base}}, \\
Q_i^{net} &= \frac{Q_i^{load} + Q_i^{cs} - Q_i^{gen} - Q_i^{sg} - Q_i^{ren} - Q_i^{pv} - Q_i^{stor}}{S_{base}}.
\end{align}
Storage discharge is therefore treated as generation, and charging stations are
additional fixed loads.

\subsection{LinDistFlow MILP}

\paragraph{Objective.}
The objective minimizes a linearized active-loss proxy:
\begin{align}
\min_x \sum_{b=1}^{m} r_b t_b.
\end{align}

\paragraph{Tree cardinality.}
\begin{align}
\sum_{b=1}^{m} \alpha_b = n-1.
\end{align}

\paragraph{Commodity-flow connectivity.}
The solver uses a single-commodity-flow spanning-tree formulation.  For the root
bus,
\begin{align}
\sum_{b\in\delta^{-}(root)} f_b - \sum_{b\in\delta^{+}(root)} f_b = -(n-1),
\end{align}
and for every non-root bus,
\begin{align}
\sum_{b\in\delta^{-}(i)} f_b - \sum_{b\in\delta^{+}(i)} f_b = 1.
\end{align}
Capacity of \(f_b\) is tied to \(\alpha_b\) using big-\(M\) logic.

\paragraph{Power-balance constraints.}
For each non-root bus,
\begin{align}
\sum_{b\in\delta^{-}(i)} P_b - \sum_{b\in\delta^{+}(i)} P_b &= P_i^{net}, \\
\sum_{b\in\delta^{-}(i)} Q_b - \sum_{b\in\delta^{+}(i)} Q_b &= Q_i^{net}.
\end{align}

\paragraph{Voltage-drop constraints.}
The current implementation uses the LinDistFlow relation relaxed by big-\(M\):
\begin{align}
v_j - v_i + 2r_b P_b + 2x_b Q_b &\le M(1-\alpha_b), \\
-(v_j - v_i + 2r_b P_b + 2x_b Q_b) &\le M(1-\alpha_b).
\end{align}

\paragraph{Thermal and absolute-value constraints.}
\begin{align}
|P_b| &\le P_b^{max}\alpha_b, \\
|Q_b| &\le Q_b^{max}\alpha_b, \\
P_b &\le t_b, \\
-P_b &\le t_b.
\end{align}

\subsection{Branch-and-Cut Configuration}

The ONR path solves the MILP through the native branch-and-cut engine with the
following current tuning:
\begin{itemize}
  \item cut family: MIR cuts,
  \item root cut rounds: 5,
  \item cuts per round: 20,
  \item branching: pseudocost,
  \item node selection: hybrid,
  \item LP node solver: simplex enabled,
  \item feasibility pump: enabled.
\end{itemize}

\subsection{Postprocessing and Verification}

The branch-status variables returned by branch-and-cut are snapped to a binary
assignment with exactly \(n-1\) closed branches.  The resulting switching plan
is then verified by running the full Newton PF solver on a modified \texttt{ACSystem}
where branch \texttt{in\_service} flags follow the optimal topology.

\subsection{Current Modeling Boundary}

The ONR module is already useful for radial feeder reconfiguration, but its
current mathematical scope is intentionally narrower than the main PF/OPF stack.
The main structural limitations are:
\begin{itemize}
  \item no hybrid-side injections from \texttt{VPP}, \texttt{Microgrid}, or \texttt{MobileStorage},
  \item no integer co-optimization of DER dispatch and topology in the current public API.
\end{itemize}
